% 356_SM_Teilchen_GF27_De_ch.tex
% Johann Pascher, 7. September 2026
% Standardmodell-Teilchen aus GF(27)* — Übersichtsdiagramm (Querformat A4)
\providecommand{\BM}{\textbf{[B]}}
\providecommand{\KM}{\textbf{[K]}}
\providecommand{\SM}{\textbf{[S]}}
\providecommand{\XM}{\textbf{[X]}}

\phantomsection
\addcontentsline{toc}{chapter}{%
  Standardmodell-Teilchen aus GF$(27)^*$ --- Übersichtsdiagramm}

\section*{Statuszeichen}
\begin{center}
\begin{tabular}{@{}lp{10.5cm}@{}}
\textbf{[B]} & \emph{Belegt} --- algebraisch bewiesen. \\[3pt]
\textbf{[K]} & \emph{Konfirmiert} --- numerisch geprüft. \\[3pt]
\textbf{[S]} & \emph{Spekulativ} --- offen. \\[3pt]
\textbf{[X]} & \emph{Experimentell} --- gesichert, FFGFT-Herleitung offen. \\
\end{tabular}
\end{center}

\section*{Beschreibung}

Dieses Querformat-Übersichtsblatt (Dok.~356) zeigt alle Standardmodell-Fermionen,
Eichbosonen und das Higgs-Boson, geordnet nach Generationen und
klassifiziert nach ihrer Herkunft aus der Galois-Feldstruktur
GF$(27)^* = \mathrm{GF}(3^3)^*$.

\textbf{Kernaussage:} Alle Fermion-Quantenzahlen des Standardmodells
($Q$, $Y$, Farbladung, Chiralität, Generationsstruktur) folgen algebraisch
aus GF$(27)^*$ ohne freie Parameter \BM\ (Dok.~346, 347, 348, 349).

\section{Fermionen}

\subsection*{Leptonen --- Quadratische Reste (QR)}

Geladene Leptonen ($e$, $\mu$, $\tau$) und Neutrinos ($\nu_e$, $\nu_\mu$,
$\nu_\tau$) entstammen dem QR-Sektor von GF$(27)^*$.
Die elektrische Ladung $Q \in \{0, -1\}$ folgt aus dem QR-Bit und
$N \bmod 3$ (Dok.~346 \BM).
Die Hyperladung $Y = -1$ folgt aus dem Legendre-Symbol als
NQR-Bit (Dok.~347 Satz~A \BM).

Alle drei Generationen entsprechen den drei Frobenius-Zyklen der
Nullstellen-Orbits der Ordnung 3 in GF$(27)^*$ \BM\ (Dok.~348 Satz~A).

\subsection*{Quarks --- Nicht-Quadratische Reste (NQR)}

Up-Quarks ($u$, $c$, $t$) und Down-Quarks ($d$, $s$, $b$) entstammen dem
NQR-Sektor von GF$(27)^*$.
Die elektrischen Ladungen $Q \in \{+2/3, -1/3\}$ folgen aus dem
NQR-Bit und $N \bmod 3$ ohne freie Parameter \BM\ (Dok.~346).
Die Farbladung $r/g/b$ wird durch den $Z_3$-Ladungsoperator
$N \bmod 3$ in GF$(3)$ erzwungen: drei Farben = drei Elemente \BM\
(Dok.~346 Satz~C).

\section{Eichbosonen}

\subsection*{Gluonen}

Die 8 Gluonen entstehen aus $4 \times Z_{13}$-Orbits $\times$ $2$
$Z_2$-Paritäten in GF$(27)^*$ \BM\ (Dok.~339).
Die Gluonzahl 8 ist strukturell erzwungen, nicht postuliert.

\subsection*{Elektroschwache Eichbosonen ($W^\pm$, $Z^0$, $\gamma$)}

Die elektroschwache Gruppe $Z_2$ folgt aus
$\mathrm{Gal}(\mathrm{GF}(9)/\mathrm{GF}(3)) = Z_2$ \KM\ (Dok.~336).
Die spontane Symmetriebrechung
$SU(2)_L \times U(1)_Y \to U(1)_\text{EM}$ ist aus drei
Galois-Bausteinen algebraisch erzwungen \BM\ (Dok.~355).

\subsection*{Higgs-Boson $H^0$}

In FFGFT ersetzt $\tilde{T} \cdot m = 1$ den Higgs-Mechanismus
(Dok.~019). Das Higgs-Boson ist experimentell gesichert \XM;
die FFGFT-Herleitung des Mechanismus ist offen.

\section{Chiralität und SU(2)$_L$-Dublett-Struktur}

Der Sd-Sektor (ungerade Grade in Cl$(6)$) entspricht dem
linkshändigen Dublett $SU(2)_L$ \BM\ (Dok.~349 Satz~D).
Der Su-Sektor (gerade Grade) ist rechtshändig.
Das rechtshändige Neutrino $\nu_R$ ist strukturell absent:
FFGFT sagt Dirac-Neutrinos voraus \BM/\SM\ (Dok.~349).

\section{Offene Punkte}

\begin{itemize}
  \item \textbf{Generations-Zuordnung (Reihenfolge):} Welches
    Orbit-Element entspricht welcher Generation (1/2/3)?
    Die Orbits sind gleichwertig; eine Zuordnung ist nicht
    abgeleitet \SM\ (Dok.~346).
  \item \textbf{CKM/PMNS-Winkel:} Die Struktur der Mischungsmatrizen
    ist belegt (Dok.~348 Sätze B, C \BM); die Zahlenwerte der
    Winkel erfordern weitere Arbeit (Dok.~348 Satz~D \SM).
  \item \textbf{Higgs-Herleitung:} FFGFT-Mechanismus
    $\tilde{T} \cdot m = 1$ ist experimentell gesichert \XM;
    formale Ausarbeitung offen (Dok.~019).
\end{itemize}
